\section{Model Training}

\subsection{Pre-training}

To enhance our language model's proficiency in generating formal proofs and reasoning through mathematical language, we further pre-train our base model~\citep{shao2024deepseekmath}. This refinement involved training on high-quality datasets that include both code and natural language mathematical content. We specifically focused on formal languages widely used in proof assistants, such as Lean, Isabelle, and Metamath. We designate this improved model as \thiswork-Base.

\subsection{Supervised Fine-tuning}

In this section, we explore the methodology and processes involved in the supervised fine-tuning (SFT) of \thiswork. Specifically, we augment the proof dataset from \lastwork~by adding detailed explanatory comments. This enhancement aims to improve the alignment between natural language descriptions and Lean 4 code, thereby facilitating better formal mathematical reasoning. Additionally, we incorporate intermediate tactic state information as an auxiliary prediction task to support the truncate-and-resume mechanism used in the Monte-Carlo Tree Search process. We refer to the resulting model as \thiswork-SFT.

\vspace{-0.1in}

\paragraph{Data Curation.}
We develop a comprehensive Lean 4 code completion dataset for the supervised fine-tuning. This dataset includes synthetic proof code derived from a wide range of formal theorems. These theorems are sourced from various projects, such as the standard Lean 4 math library Mathlib4 \citep{mathlib}, synthetic theorems from \lastwork~\citep{xin2024deepseek} and Lean Workbook \citep{ying2024lean}, and validation sets from the miniF2F \citep{zheng2021minif2f} and ProofNet \citep{azerbayev2023proofnet} benchmarks. To augment the formal proof data, we employed an expert iteration process \citep{polu2020generative}. This involves generating proofs using the language model, verifying the generated proof data, retraining the model with the verified data, and then using the optimized model to generate additional proof data. Between each iteration, we use DeepSeek-Coder V2 236B \citep{zhu2024deepseek} to annotate the thought process before the proof code as comments. Finally, we tailor these data for the truncate-and-resume mechanism for Monte-Carlo Tree Search (details in Section~\ref{sec:tree-abstraction}). The resulting proof dataset consists of 9,645k sequences.

\vspace{-0.1in}

\paragraph{Thought-augmented Proof Generation. } \label{sec:thought-augmented_proof_generation}
In \lastwork, we identified a significant gap between problem-solving strategies in natural language and theorem proving in Lean. In natural language, models generate detailed deduction steps to construct proofs, whereas in Lean, they often rely on a sequence of high-level tactic calls to brute-force solutions. These high-level tactics, while effective, obscure their internal workings and outcomes, hindering the model's ability to resolve complex proof goals with structured mathematical reasoning. To address this issue, we develop an approach that incorporates natural language reasoning before generating theorem proof code. Similar to Lean-STaR \citep{lin2024lean}, which performs isolated chain-of-thought reasoning \citep{wei2022chain, feng2024towards} before each proof step, our method integrates this reasoning directly as comments within the proof code. We use the DeepSeek-Coder V2 236B \citep{zhu2024deepseek} to enhance existing data in \lastwork~in two ways: first, by inserting a complete natural language solution at the beginning of the proof block, and second, by alternately inserting specific natural language steps for corresponding Lean tactics. Training the model with this data format enforces it to propose complete mathematical reasoning at the beginning of the proof block and detailed step planning before each tactic. This approach successfully develops new behaviors, employing delicate mathematical thinking to guide the generation of tactics. In the training data, two distinct guiding prompts are used to differentiate between the CoT (Chain of Thought) mode and the non-CoT mode for proof code completion. Examples of input and output in both modes can be found in Appendix~\ref{app:promt_examples}.

\vspace{-0.1in}

\paragraph{Prompt Augmentation with Tactic State Information. } \label{sec:sft-tactic-state}
To implement the truncate-and-resume mechanism for Monte-Carlo Tree Search, we needed to extract tactic information from the code generated by the model. We enhanced the Lean REPL \citep[Read-Eval-Print Loop;][]{repl} with data extraction tools from the LeanDojo \citep{yang2024leandojo} project. This allowed us to extract tactic information in triples, which include the position of each tactic, as well as the tactic states before and after its application. This information helps us identify the specific tactic code that triggers verification errors (used in the expansion step for tree search, see Section~\ref{para:expansion}). For each tactic in a generated valid formal proof, we insert the tactic state returned by the verifier as a comment "/- tactic state: ... -/". During training, we use all tokens following "/- tactic state: " as responses to calculate the supervised fine-tuning loss, while the tokens before this comment is used as prompts and do not contribute to the training loss calculation.

\vspace{-0.1in}

\paragraph{Training Setting. }
We conduct supervised fine-tuning based on the pre-trained model and train for 9B tokens, using a batch size of 2,048 and a constant learning rate of 1e-4. The training process begins with 100 warm-up steps to stabilize the learning dynamics. Training examples are randomly concatenated to form sequences, with a maximum context length of 4,096 tokens.

\subsection{Reinforcement Learning from Proof Assistant Feedback}

Reinforcement learning (RL) has been proven effective in enhancing the mathematical reasoning capabilities of supervised fine-tuned language models~\citep{shao2024deepseekmath}. To further advance \thiswork-SFT, we incorporate a reinforcement learning phase, resulting in the model \thiswork-RL. This phase leverages RL to enhance performance based on verification feedback from the Lean 4 prover. The specifics of this RL process are detailed below.

\paragraph{Prompts.}
In the reinforcement learning stage, we use a subset of theorem statements from the supervised fine-tuning dataset as training prompts. We select theorems for which \thiswork-SFT has a moderate success rate in generating correct proofs upon multiple attempts. This ensures that the model has room for improvement while still being able to receive positive feedback. After filtering, we retain approximately 4.5k unique theorem statements. Each theorem is prefixed with both CoT and non-CoT guiding prompts to enhance the model's proof generation capabilities in both modes.

\paragraph{Rewards.}
When training LLMs via RL, a trained reward model typically provides feedback signals. In contrast, formal theorem proving benefits from the rigorous verification of generated proofs by proof assistants, offering a significant advantage. Specifically, each generated proof receives a reward of 1 if verified as correct, and 0 otherwise. While this binary reward signal is accurate, it is also sparse, especially for theorems that are challenging for the supervised fine-tuned model. To mitigate this sparsity, we select training prompts that are challenging yet achievable for the supervised fine-tuned model, as described above.

\paragraph{Reinforcement Learning Algorithm.}
We employ the Group Relative Policy Optimization \citep[GRPO;][]{shao2024deepseekmath} as our RL algorithm, which has demonstrated superior effectiveness and efficiency compared to PPO \citep{schulman2017proximal}, primarily because it eliminates the necessity of training an additional critic model. Specifically, GRPO samples a group of candidate proofs for each theorem prompt and optimizes the model based on the relative rewards of the outputs within the group. Our prompt selection strategy is designed to likely include both correct and incorrect proofs among the candidates, aligning well with the group-relative nature of GRPO and thereby enhancing the training process.

\paragraph{Training Setting. }
We conduct RL training based on the SFT model, which serves as both the initial model and the reference model for imposing the Kullback-Leibler (KL) divergence penalty. We use a constant learning rate of 5e-6, and the KL penalty coefficient is set to 0.02. For each theorem, we sample a group of 32 candidate proofs, with maximum length set to 2,048. The training batch size is configured to 512.

\subsection{Evaluation} \label{sec:sft-eval}

\paragraph{Benchmarks.}
We evaluate theorem-proving performance on the following benchmarks to compare model capabilities after each training stage:
\begin{itemize}
    \item \textbf{MiniF2F}~\citep{zheng2021minif2f} focuses on formal problem-solving skills for high-school level exercises and competitions, such as AMC, AIME, and IMO, with an emphasis on algebra and number theory. The benchmark includes 244 validation and 244 test problems, originally in Lean 3 and manually converted to Lean 4.9.0, based on the version provided by \citet{minif2f_solution}.
    \item \textbf{ProofNet}~\citep{azerbayev2023proofnet} evaluates formal theorem-proving capabilities at the undergraduate level in mathematics. It comprises 185 validation and 186 test problems from widely-used undergraduate textbooks, covering real and complex analysis, linear algebra, abstract algebra, and topology. These problems were initially in Lean 3 and manually converted to Lean 4.9.0.
\end{itemize}

\begin{figure}[t]
  \begin{minipage}{0.55\textwidth}
    \centering
    \includegraphics[width=0.98\textwidth]{figures/sample128.pdf}
  \end{minipage}
  \begin{minipage}{0.45\textwidth}
    \centering
    \footnotesize
    \begin{tabular}{lcc}
    \toprule % \vspace{0.02in}
        \multirow{2}{*}{Model} & \multicolumn{2}{c}{Pass@128} \\
        & miniF2F-test & ProofNet-test \\
        \toprule
        Base (3-shot) & $29.7\% \pm 0.5\%$ & $9.7\% \pm 0.7\%$ \\
        \midrule
        SFT (non-CoT) & $49.8\% \pm 0.3\%$ & $15.9\% \pm 0.5\%$ \\
        SFT (CoT) & $50.4\% \pm 0.4\%$ & $15.9\% \pm 0.6\%$ \\
        \midrule
        RL (non-CoT) & $50.5\% \pm 0.6\%$ & $17.5\% \pm 0.5\%$\\
        RL (CoT) & $51.6\% \pm 0.5\%$ & $18.2\% \pm 0.5\%$ \\
        \bottomrule
    \end{tabular}
  \end{minipage}
  \caption{\textbf{Comparison of model capabilities at different training stages.} "CoT" and "non-CoT" refer to evaluations using two guiding prompts. The shaded region represents the range of standard deviations around the mean values. The notation $\mu\pm\sigma$ indicates the average accuracy $\mu$ and the standard deviation $\sigma$.}
  \label{fig:sample128}
\end{figure}

\paragraph{Prompting Configurations.}
For each proof attempt of \thiswork-Base, we independently sample three proof demonstrations from the validation set to construct the few-shot prompts. For the miniF2F benchmark, we use human-written proofs from \citet{minif2f_solution}, while for the ProofNet benchmark, we use correct proofs generated by \thiswork-RL as few-shot demonstrations. For \thiswork-SFT and \thiswork-RL, we employ two types of guiding prompts: one that encourages chain-of-thought (CoT) reasoning before each proof step, and one that does not (non-CoT). Detailed examples are provided in Appendix~\ref{app:promt_examples}.

\paragraph{Metric.}
We evaluate theorem-proving performance using the pass@$K$ accuracy metric, which measures the model's success in generating a correct proof within $K$ attempts. Each model is deployed on a single A100-40G GPU, utilizing the vLLM framework~\citep{kwon2023efficient} for sample generation. The sampling parameters are set with a temperature of 1, a top-p value of 0.95, and a maximum token limit of 2,048. The generated proofs are then verified using the Lean 4 theorem prover. For this verification, we import Mathlib4 \citep{mathlib} and Aesop \citep{limperg2023aesop} to access predefined premises and tactics. The verification process is subject to a time limit of 300 seconds.

\paragraph{Comparison across Training Stages.}
Figure~\ref{fig:sample128} presents a comparative analysis of each training stage on the miniF2F and ProofNet datasets. Our base model, \thiswork-Base, achieves a notable pass rate, solving nearly one-third of the problems on the test set of the miniF2F benchmark using 3-shot prompting. The supervised fine-tuning stage, resulting in \thiswork-SFT, significantly outperforms the base model, with Pass@128 accuracy increasing by approximately two-thirds on miniF2F and doubling on ProofNet. The subsequent reinforcement learning stage further enhances the model's performance, improving Pass@$K$ accuracy across all values of $K$. In contrast to findings in natural language mathematics, such as those reported in DeepSeekMath~\citep{shao2024deepseekmath}, where reinforcement learning primarily boosts the correct response from TopK, we observe a genuine enhancement of fundamental capabilities in formal theorem proving. This improvement is evident not only with a small sample budget but also remains stable as the sample budget increases. This conclusion is further supported by later Monte-Carlo Tree Search experiments with larger sample budgets, as discussed in Section~\ref{sec:large-train-eval}.

\paragraph{Comparison between CoT and non-CoT.}
We compare the performance of non-CoT and CoT generation modes for both \thiswork-SFT and \thiswork-RL. The results in Figure~\ref{fig:sample128} demonstrate that the CoT mode consistently outperforms the non-CoT mode across most settings. Specifically, \thiswork-RL, leveraging these enhanced theorem-proving patterns, achieves superior performance on both benchmarks, with an average accuracy of $51.6\%$ on miniF2F and $18.2\%$ on ProofNet. The integration of natural language reasoning in CoT mode significantly enhances the planning and execution of formal proof writing. For a detailed comparison of proof strategies with and without the use of natural language chain-of-thought, refer to the examples provided in Appendix~\ref{app:promt_examples}.
